\begin{proof}[Proof of \cref{obj:thm:support-transition}]
Throughout, fix \(p=3/10\) and \(\varepsilon=\log 3\). Set
\[
A_5=\{S_3,S_4,S_6,S_8,S_9\},
\qquad
A_3=\{S_1,S_6,S_8\},
\qquad
\iota:\mathbb R^2\longrightarrow\mathbb R^4,
\quad
\iota(\theta)=\left(\frac3{10},\mu_0,\mu_1,\log 3\right),
\]
where \(\theta=(\mu_0,\mu_1)^\top\). We use the objective, rays, and certificates of
\cref{obj:synth_4,obj:synth_5}, with the pattern enumeration in
\cref{obj:synth_8}.

\begin{enumerate}
\item \emph{Construct the open slices.}
By \cref{obj:thm:five-output-region,obj:thm:three-output-region},
\[
\mathcal R_5,\mathcal R_3\ \text{are open},\qquad
\iota\bigl((1/2,1/2)^\top\bigr)\in\mathcal R_5,
\qquad
\iota\bigl((1/10,1/10)^\top\bigr)\in\mathcal R_3.
\]
Define
\[
\mathcal V_5
=\iota^{-1}(\mathcal R_5)\cap\{\theta:\mu_0>3/10\},
\qquad
\mathcal V_3
=\iota^{-1}(\mathcal R_3)\cap\{\theta:\mu_0<3/10\}.
\]
The map \(\iota\) is continuous, so both sets are open. Their respective centers belong to them because
\[
\frac1{10}<\frac3{10}<\frac12.
\]
They are disjoint: membership in their intersection would require both
\(\mu_0>3/10\) and \(\mu_0<3/10\). Moreover,
\[
\theta\in\mathcal V_5\Longrightarrow\iota(\theta)\in\mathcal R_5,
\qquad
\theta\in\mathcal V_3\Longrightarrow\iota(\theta)\in\mathcal R_3.
\]
% lean: CausalSmith.Stat.LdpAteEfficiencySurface.supportRegion_slices

\item \emph{Preserve the three distinct scores near the low-mean point.}
For this neighborhood construction, let
\[
\mathcal O=(0,1)^3\times(0,\infty),\qquad
\xi=(p_\xi,\mu_{0,\xi},\mu_{1,\xi},\varepsilon_\xi)\in\mathcal O,
\qquad
\theta_\xi=(\mu_{0,\xi},\mu_{1,\xi})^\top.
\]
Write the record probabilities, rays, masses, and gradients at these parameters as
\[
\begin{aligned}
\pi_\xi&=\bigl((1-p_\xi)(1-\mu_{0,\xi}),
 (1-p_\xi)\mu_{0,\xi},
 p_\xi(1-\mu_{1,\xi}),p_\xi\mu_{1,\xi}\bigr)^\top,\\
b_S(\xi)&=\mathbf1_4+(e^{\varepsilon_\xi}-1)\mathbf1_S,
\qquad h_S(\xi)=\pi_\xi^\top b_S(\xi),\\
g_S(\xi)&=(e^{\varepsilon_\xi}-1)
\begin{pmatrix}
(1-p_\xi)\{\mathbf1_S(1)-\mathbf1_S(0)\}\\
p_\xi\{\mathbf1_S(3)-\mathbf1_S(2)\}
\end{pmatrix},
\qquad
\mathbf1_2=\begin{pmatrix}1\\1\end{pmatrix}.
\end{aligned}
\]
Here \(\mathbf1_S\) is the indicator vector of the pattern \(S\). In particular,
\(h_S(\xi)>0\) for every \(S\).

Define the canonical three-pattern weights and their stationarity coefficients by
\[
\alpha^{\mathrm{can}}_S(\xi)=
\begin{cases}
(e^{\varepsilon_\xi}+2)^{-1},&S\in A_3,\\
0,&S\notin A_3,
\end{cases}
\]
\[
a_3(\xi)=\frac{2}{e^{\varepsilon_\xi}+2}
 \sum_{S\in A_3}\frac{(g_S(\xi)^\top\mathbf1_2)^2}{h_S(\xi)},
\qquad
\ell_3(\xi)=\frac{2}{e^{\varepsilon_\xi}+2}
 \sum_{S\in A_3}
 \frac{(g_S(\xi)^\top v(0))(g_S(\xi)^\top\mathbf1_2)}{h_S(\xi)}.
\]
Expansion of the quadratic objective gives
\[
\frac{\partial}{\partial t}
F_{\theta_\xi}\bigl(\alpha^{\mathrm{can}}(\xi),t\bigr)
=\ell_3(\xi)+a_3(\xi)t,
\]
where the objective is evaluated at \(p_\xi,\varepsilon_\xi\). The \(S_1\) summand yields
\[
a_3(\xi)\ge
\frac{2(e^{\varepsilon_\xi}-1)^2(1-p_\xi)^2}
 {(e^{\varepsilon_\xi}+2)h_{S_1}(\xi)}>0.
\]
Thus the unique stationary coordinate is the function
\[
t_3:\mathcal O\longrightarrow\mathbb R,
\qquad
t_3(\xi)=-\frac{\ell_3(\xi)}{a_3(\xi)}.
\]
All denominators in these formulas are positive, so \(t_3\) and the three functions
\[
r_{3,S}(\xi)=
\frac{g_S(\xi)^\top v(t_3(\xi))}{h_S(\xi)},
\qquad S\in A_3,
\]
are continuous at
\[
\xi_{\mathrm{low}}
=\left(\frac3{10},\frac1{10},\frac1{10},\log 3\right).
\]

To calculate their values, the record probabilities at this point are
\[
\pi_{\xi_{\mathrm{low}}}
=\left(\frac{63}{100},\frac7{100},\frac{27}{100},\frac3{100}\right)^\top.
\]
Using \(S_1=\{0\}\), \(S_6=\{1,2\}\), and \(S_8=\{3\}\), the definitions give
\[
\begin{array}{c|c|c}
S&h_S(\xi_{\mathrm{low}})&g_S(\xi_{\mathrm{low}})\\ \hline
S_1&113/50&(-7/5,0)^\top\\
S_6&42/25&(7/5,-3/5)^\top\\
S_8&53/50&(0,3/5)^\top
\end{array}
\]
and hence
\[
\ell_3(\xi_{\mathrm{low}})=\frac8{371},
\qquad
a_3(\xi_{\mathrm{low}})=\frac{79880}{125769},
\qquad
t_3(\xi_{\mathrm{low}})=-\frac{339}{9985}.
\]
Substitution into the projected-score formula gives
\[
r_{3,S_1}(\xi_{\mathrm{low}})=\frac{42}{1997},
\qquad
r_{3,S_6}(\xi_{\mathrm{low}})=-\frac{1491}{3994},
\qquad
r_{3,S_8}(\xi_{\mathrm{low}})=\frac{1092}{1997}.
\]
These values are pairwise distinct: the middle value is negative, the others are positive, and
\[
\frac{1092}{1997}-\frac{42}{1997}=\frac{1050}{1997}>0.
\]
For each ordered pair \(S,T\in A_3\) with \(S\ne T\), continuity of their difference therefore supplies an open neighborhood \(D_{S,T}\) satisfying
\[
\xi_{\mathrm{low}}\in D_{S,T}
\subseteq
\{\xi\in\mathcal O:r_{3,S}(\xi)-r_{3,T}(\xi)\ne0\}.
\]
Define
\[
D_3=\bigcap_{\substack{S,T\in A_3\\S\ne T}}D_{S,T}.
\]
This finite intersection is open, contains \(\xi_{\mathrm{low}}\), and satisfies
\[
\xi\in D_3,\quad S,T\in A_3,\quad S\ne T
\Longrightarrow r_{3,S}(\xi)\ne r_{3,T}(\xi).
\]
% lean: CausalSmith.Stat.LdpAteEfficiencySurface.exists_lowMean_r3_distinctScore_neighborhood

\item \emph{Choose the final neighborhoods.}
Set
\[
U_5=\mathcal V_5,
\qquad
U_3=\mathcal V_3\cap\iota^{-1}(D_3).
\]
Continuity of \(\iota\) makes \(U_3\) open. The center conditions established above give
\[
(1/2,1/2)^\top\in U_5,
\qquad
(1/10,1/10)^\top\in U_3.
\]
Since \(U_3\subseteq\mathcal V_3\), the sets \(U_5,U_3\) are disjoint. Every point in either neighborhood has interior arm means, by its corresponding region membership.
% lean: CausalSmith.Stat.LdpAteEfficiencySurface.support_transition

\item \emph{Obtain the five-pattern optimum and support rigidity.}
Fix \(\theta\in U_5\). Membership in \(\mathcal R_5\), together with
\cref{obj:def:support-regions,obj:synth_5}, supplies
\[
\alpha^{(5)}\in\mathcal A_\varepsilon,\qquad
t_{\mathrm{cert},5}\in\mathbb R,\qquad
\eta^{(5)}\in\mathbb R^4,
\]
with
\[
\alpha^{(5)}_S>0\quad(S\in A_5),\qquad
\alpha^{(5)}_S=0\quad(S\notin A_5),
\qquad
\frac{\partial}{\partial t}
F_\theta(\alpha^{(5)},t_{\mathrm{cert},5})=0,
\]
and
\[
f_S(t_{\mathrm{cert},5})=
 b_S^\top\eta^{(5)}\quad(S\in A_5),
\qquad
f_S(t_{\mathrm{cert},5})<
 b_S^\top\eta^{(5)}\quad(S\notin A_5).
\]
Here \(f_S(t)=(g_S^\top v(t))^2/h_S(\theta)\), as in
\cref{obj:synth_5}. Stationarity gives the exact expansion
\[
F_\theta(\alpha^{(5)},t)-F_\theta(\alpha^{(5)},t_{\mathrm{cert},5})
=(t-t_{\mathrm{cert},5})^2
 \sum_{S\in\mathcal S}
 \alpha^{(5)}_S\frac{(g_S^\top\mathbf1_2)^2}{h_S(\theta)}
\ge0.
\]
For every feasible \(\beta\), normalization and the dual inequalities give
\[
\begin{aligned}
\inf_{t\in\mathbb R}F_\theta(\beta,t)
&\le F_\theta(\beta,t_{\mathrm{cert},5})\\
&\le\sum_{S\in\mathcal S}\beta_S b_S^\top\eta^{(5)}
=\mathbf1_4^\top\eta^{(5)}
=F_\theta(\alpha^{(5)},t_{\mathrm{cert},5}).
\end{aligned}
\]
The last equality uses the active dual equalities and the vanishing inactive weights. Since \(\alpha^{(5)}\) attains its minimum at the certifying coordinate,
\[
J^*(\theta,p,\varepsilon)
=\inf_tF_\theta(\alpha^{(5)},t)
=F_\theta(\alpha^{(5)},t_{\mathrm{cert},5})
=\mathbf1_4^\top\eta^{(5)}.
\]

If a feasible \(\beta\) satisfies \(J^*=\inf_tF_\theta(\beta,t)\), the preceding inequalities are equalities. Consequently,
\[
0=\sum_{S\in\mathcal S}\beta_S
 \bigl(b_S^\top\eta^{(5)}-f_S(t_{\mathrm{cert},5})\bigr).
\]
Every summand is nonnegative, and every inactive slack is strictly positive. Therefore
\[
\beta_S=0\qquad(S\notin A_5).
\]
% lean: CausalSmith.Stat.LdpAteEfficiencySurface.strictCertificate_profile_optimal_and_rigid

\item \emph{Determine the five-pattern mixture uniquely.}
For a feasible \(\beta\) vanishing outside \(A_5\), the four normalization equations, using \(e^\varepsilon=3\), are
\[
\begin{aligned}
3\beta_{S_3}+\beta_{S_4}+\beta_{S_6}+\beta_{S_8}+3\beta_{S_9}&=1,\\
3\beta_{S_3}+\beta_{S_4}+3\beta_{S_6}+\beta_{S_8}+\beta_{S_9}&=1,\\
\beta_{S_3}+3\beta_{S_4}+3\beta_{S_6}+\beta_{S_8}+\beta_{S_9}&=1,\\
\beta_{S_3}+\beta_{S_4}+\beta_{S_6}+3\beta_{S_8}+3\beta_{S_9}&=1.
\end{aligned}
\]
Subtracting these equations yields
\[
\beta_{S_3}=\beta_{S_4}=\beta_{S_8},
\qquad
\beta_{S_6}=\beta_{S_9},
\qquad
5\beta_{S_3}+4\beta_{S_6}=1.
\]
Thus, with the mixture coordinate defined by
\[
\lambda(\beta)=4\beta_{S_6},
\]
the weights satisfy
\[
\beta_{S_6}=\beta_{S_9}=\frac{\lambda(\beta)}4,
\qquad
\beta_{S_3}=\beta_{S_4}=\beta_{S_8}
=\frac{1-\lambda(\beta)}5.
\]

Define the two normalized component objectives and their coefficients by
\[
B_5(t)=\frac{f_{S_6}(t)+f_{S_9}(t)}4,
\qquad
T_5(t)=\frac{f_{S_3}(t)+f_{S_4}(t)+f_{S_8}(t)}5,
\]
\[
C_{\mathrm{bin}}=\frac1{100}
 \left(\frac1{h_{S_6}(\theta)}+\frac1{h_{S_9}(\theta)}\right)>0,
\qquad
C_{\mathrm{tri}}=\frac9{125}
 \left(\frac1{h_{S_4}(\theta)}+\frac1{h_{S_8}(\theta)}\right)>0.
\]
The pattern gradients are
\[
g_{S_3}=0,\quad
g_{S_4}=(0,-3/5)^\top,\quad
g_{S_6}=(7/5,-3/5)^\top,\quad
g_{S_8}=(0,3/5)^\top,\quad
g_{S_9}=(-7/5,3/5)^\top,
\]
so direct substitution gives
\[
B_5(t)=C_{\mathrm{bin}}(4t-3)^2,
\qquad
T_5(t)=C_{\mathrm{tri}}(t+1)^2.
\]
Also,
\[
\frac{b_{S_6}+b_{S_9}}4
=\frac{b_{S_3}+b_{S_4}+b_{S_8}}5
=\mathbf1_4.
\]
The active dual equalities therefore imply
\[
B_5(t_{\mathrm{cert},5})
=T_5(t_{\mathrm{cert},5})
=\mathbf1_4^\top\eta^{(5)}.
\]
In particular,
\[
T_5'(t_{\mathrm{cert},5})
=2C_{\mathrm{tri}}(t_{\mathrm{cert},5}+1)\ne0.
\]
Indeed, a zero derivative would give \(t_{\mathrm{cert},5}=-1\), whereas the component equality would then require
\(49C_{\mathrm{bin}}=0\), contradicting \(C_{\mathrm{bin}}>0\).

For the certificate weights, the mixture representation gives
\[
F_\theta(\alpha^{(5)},t)
=\lambda(\alpha^{(5)})B_5(t)
 +(1-\lambda(\alpha^{(5)}))T_5(t).
\]
Its stationarity consequently implies
\[
\lambda(\alpha^{(5)})B_5'(t_{\mathrm{cert},5})
 +(1-\lambda(\alpha^{(5)}))T_5'(t_{\mathrm{cert},5})=0.
\]
The derivative difference is nonzero: if the two derivatives were equal, this equation would force \(T_5'(t_{\mathrm{cert},5})=0\). Solving gives
\[
\lambda(\alpha^{(5)})
=-\frac{T_5'(t_{\mathrm{cert},5})}
 {B_5'(t_{\mathrm{cert},5})-T_5'(t_{\mathrm{cert},5})}.
\]
For any feasible optimal \(\beta\), the preceding support-rigidity argument applies, and
\[
F_\theta(\beta,t_{\mathrm{cert},5})
=J^*=\inf_tF_\theta(\beta,t).
\]
Hence \(\beta\) is stationary at the same coordinate and satisfies the same mixture equation. Thus
\[
\lambda(\beta)=\lambda(\alpha^{(5)}),
\qquad
\beta=\alpha^{(5)}.
\]
This proves the required uniqueness among all feasible optimal weights.
% lean: CausalSmith.Stat.LdpAteEfficiencySurface.uniqueOptimalSupport_of_mem_R5

\item \emph{Establish exact cardinality five.}
We first record positivity of the optimized information throughout the interior fixed-parameter slice. Define the feasible singleton design by
\[
\alpha^{\mathrm{rr}}_S=
\begin{cases}
1/6,&S\in\{S_1,S_2,S_4,S_8\},\\
0,&\text{otherwise}.
\end{cases}
\]
Each input row sums to \((3+1+1+1)/6=1\). Its objective satisfies, for every \(t\in\mathbb R\),
\[
F_\theta(\alpha^{\mathrm{rr}},t)
\ge\frac{49}{150h_{S_1}(\theta)}t^2,
\qquad
F_\theta(\alpha^{\mathrm{rr}},t)
\ge\frac{3}{50h_{S_8}(\theta)}(t+1)^2.
\]
It is a coercive quadratic, so choose
\[
t_{\mathrm{rr}}\in\arg\min_{t\in\mathbb R}
F_\theta(\alpha^{\mathrm{rr}},t).
\]
A zero minimum would force both \(t_{\mathrm{rr}}=0\) and
\(t_{\mathrm{rr}}=-1\). Therefore
\[
0<F_\theta(\alpha^{\mathrm{rr}},t_{\mathrm{rr}})
\le J^*(\theta,p,\varepsilon).
\]
% lean: CausalSmith.Stat.LdpAteEfficiencySurface.Jstar_pos_interior

For the five-pattern certificate, define
\[
I_5=I_\theta(\alpha^{(5)}),
\qquad
c=\begin{pmatrix}-1\\1\end{pmatrix},
\qquad
w_5=\frac{v(t_{\mathrm{cert},5})}{J^*(\theta,p,\varepsilon)}.
\]
The matrix \(I_5\) is positive semidefinite, since it is a sum of nonnegative multiples of \(g_Sg_S^\top\). Stationarity and the certificate value give
\[
\mathbf1_2^\top I_5v(t_{\mathrm{cert},5})=0,
\qquad
v(t_{\mathrm{cert},5})^\top I_5v(t_{\mathrm{cert},5})=J^*.
\]
The first identity makes the two coordinates of \(I_5v(t_{\mathrm{cert},5})\) opposite; the second makes its second coordinate \(J^*\). Hence
\[
I_5v(t_{\mathrm{cert},5})=J^*c,
\qquad
I_5w_5=c,
\qquad
c^\top w_5=\frac1{J^*}=V^*.
\]
Every other solution \(w_{\mathrm{alt}}\in\mathbb R^2\) of \(I_5w_{\mathrm{alt}}=c\) has the same contrast value, because
\[
c^\top(w_{\mathrm{alt}}-w_5)
=w_5^\top I_5(w_{\mathrm{alt}}-w_5)=0.
\]
The contrast-variance convention in \cref{obj:thm:five-output-region} therefore gives variance \(V^*>0\), equal to its nonnegative part.

For the associated staircase channel,
\[
(P_\theta Q_{\alpha^{(5)}})(\{S\})
=\alpha^{(5)}_S h_S(\theta).
\]
The positive masses \(h_S(\theta)\) make precisely the five patterns in \(A_5\) active, so
\[
\kappa(Q_{\alpha^{(5)}})=5.
\]
Feasibility supplies unit row sums, and the ray entries \(1,3\) supply the privacy inequalities. Finally, \cref{obj:thm:five-output-region} gives
\[
\kappa(Q)\ge5
\]
for every attaining stationary \(\varepsilon\)-LDP channel on any measurable output space, since \(\iota(\theta)\in\mathcal R_5\).
% lean: CausalSmith.Stat.LdpAteEfficiencySurface.optimalCardinality_of_mem_R5

\item \emph{Obtain the unique three-pattern optimum.}
Fix \(\theta\in U_3\). Its membership in \(\mathcal R_3\) supplies certificate data
\[
\alpha^{(3)}\in\mathcal A_\varepsilon,\qquad
t_{\mathrm{cert},3}\in\mathbb R,\qquad
\eta^{(3)}\in\mathbb R^4,
\]
with support \(A_3\), positive active weights, stationarity, active dual equalities, and strict inactive dual inequalities. The quadratic expansion and normalization calculation used above give
\[
J^*
=\inf_tF_\theta(\alpha^{(3)},t)
=F_\theta(\alpha^{(3)},t_{\mathrm{cert},3})
=\mathbf1_4^\top\eta^{(3)},
\]
and, for every feasible \(\beta\),
\[
F_\theta(\beta,t_{\mathrm{cert},3})
\le\mathbf1_4^\top\eta^{(3)}=J^*.
\]
If \(J^*=\inf_tF_\theta(\beta,t)\), equality follows and
\[
0=\sum_{S\in\mathcal S}\beta_S
 \bigl(b_S^\top\eta^{(3)}-f_S(t_{\mathrm{cert},3})\bigr),
\qquad
\beta_S=0\quad(S\notin A_3).
\]

For any feasible weights supported on \(A_3\), the normalization rows are
\[
\begin{aligned}
3\beta_{S_1}+\beta_{S_6}+\beta_{S_8}&=1,\\
\beta_{S_1}+3\beta_{S_6}+\beta_{S_8}&=1,\\
\beta_{S_1}+\beta_{S_6}+3\beta_{S_8}&=1,
\end{aligned}
\]
with the second row also applying to the other record in \(S_6\). Subtraction and substitution give
\[
\beta_{S_1}=\beta_{S_6}=\beta_{S_8}=\frac15.
\]
These equations apply both to \(\alpha^{(3)}\) and to every feasible optimal \(\beta\). Consequently,
\[
\alpha^{(3)}_S=
\begin{cases}
1/5,&S\in A_3,\\
0,&S\notin A_3,
\end{cases}
\qquad
\beta=\alpha^{(3)}.
\]
% lean: CausalSmith.Stat.LdpAteEfficiencySurface.uniqueOptimalSupport_of_mem_R3

\item \emph{Construct an attaining channel with three active outputs.}
The canonical weights just obtained agree with
\(\alpha^{\mathrm{can}}(\iota(\theta))\). Their stationarity equation is
\[
0=\ell_3(\iota(\theta))+a_3(\iota(\theta))t_{\mathrm{cert},3}.
\]
Since \(a_3(\iota(\theta))>0\),
\[
t_{\mathrm{cert},3}=t_3(\iota(\theta)).
\]
The information matrix is positive semidefinite, and its certifying minimum equals \(J^*>0\). The same two-coordinate calculation therefore gives
\[
I_\theta(\alpha^{(3)})v(t_{\mathrm{cert},3})=J^*c,
\qquad
I_\theta(\alpha^{(3)})
 \frac{v(t_{\mathrm{cert},3})}{J^*}=c,
\qquad
c^\top\frac{v(t_{\mathrm{cert},3})}{J^*}=V^*>0.
\]
Thus the feasible staircase channel attains the required contrast variance. Its output probabilities satisfy
\[
(P_\theta Q_{\alpha^{(3)}})(\{S\})
=\alpha^{(3)}_S h_S(\theta),
\qquad
\kappa(Q_{\alpha^{(3)}})=|A_3|=3.
\]
% lean: CausalSmith.Stat.LdpAteEfficiencySurface.r3Certificate_staircase_attains_three

\item \emph{Prove the universal three-output lower bound.}
Let \(Q\) be any attaining stationary \(\varepsilon\)-LDP channel on a measurable output space \(Z\), and define
\[
I_Q=I_\theta(Q).
\]
Its contrast variance is the finite positive value \(V^*=1/J^*\). Hence the contrast is in the range of \(I_Q\), and we may choose
\[
w_Q\in\mathbb R^2,\qquad
I_Qw_Q=c,\qquad
c^\top w_Q=\frac1{J^*}.
\]
Positive semidefiniteness and Cauchy--Schwarz give
\[
1=(v(t_{\mathrm{cert},3})^\top I_Qw_Q)^2
\le
\bigl(v(t_{\mathrm{cert},3})^\top I_Qv(t_{\mathrm{cert},3})\bigr)
 \bigl(w_Q^\top I_Qw_Q\bigr).
\]
Since \(w_Q^\top I_Qw_Q=c^\top w_Q=1/J^*\), this implies
\[
J^*\le
v(t_{\mathrm{cert},3})^\top I_Qv(t_{\mathrm{cert},3}).
\]

Apply \cref{obj:lem:staircase-refinement}, using the interior parameters and the constant positive privacy budget. It supplies feasible weights \(\bar\alpha\) and a probability kernel \(K\) with
\[
Q=K\circ Q_{\bar\alpha},
\qquad
I_\theta(\bar\alpha)-I_Q\succeq0.
\]
The directional lower bound, refinement inequality, and certificate upper bound form the chain
\[
J^*
\le v(t_{\mathrm{cert},3})^\top I_Qv(t_{\mathrm{cert},3})
\le F_\theta(\bar\alpha,t_{\mathrm{cert},3})
\le J^*.
\]
Thus every inequality is an equality. Equality in the dual bound gives
\[
0=\sum_{S\in\mathcal S}\bar\alpha_S
 \bigl(b_S^\top\eta^{(3)}-f_S(t_{\mathrm{cert},3})\bigr),
\qquad
\bar\alpha_S=0\quad(S\notin A_3).
\]
The normalization equations above then force
\[
\bar\alpha_{S_1}=\bar\alpha_{S_6}=\bar\alpha_{S_8}=\frac15.
\]

Because \(\iota(\theta)\in D_3\) and
\(t_{\mathrm{cert},3}=t_3(\iota(\theta))\), the active projected scores satisfy
\[
\rho_S(\theta,t_{\mathrm{cert},3})
\ne\rho_T(\theta,t_{\mathrm{cert},3})
\qquad(S,T\in A_3,\ S\ne T).
\]
Equality in the refinement direction, together with
\cref{obj:lem:staircase-refinement}, consequently implies
\[
K(S,\cdot)\perp K(T,\cdot)
\qquad(S,T\in A_3,\ S\ne T).
\]

Define the output law by
\[
\nu_Q=P_\theta Q
=\sum_{S\in A_3}\frac{h_S(\theta)}5K(S,\cdot).
\]
In particular, for every measurable output event \(E\) and every \(S\in A_3\),
\[
\nu_Q(E)\ge\frac{h_S(\theta)}5K(S,E),
\qquad
K(S,\cdot)\ll\nu_Q.
\]
If \(Z\) is infinite, the cardinality convention gives
\(\kappa(Q)=+\infty\ge3\).

If \(Z\) is finite, each probability measure \(K(S,\cdot)\), for \(S\in A_3\), satisfies
\[
1=K(S,Z)\le\sum_{z\in Z}K(S,\{z\}).
\]
Thus choose, for every \(S\in A_3\), a point
\[
z_S\in Z,\qquad K(S,\{z_S\})>0.
\]
Absolute continuity transfers positivity to the output law: indeed,
\[
\nu_Q(\{z_S\})=0\Longrightarrow K(S,\{z_S\})=0,
\]
so \(\nu_Q(\{z_S\})>0\).

To prove that these points are distinct, fix distinct \(S,T\in A_3\). Mutual singularity supplies a measurable set \(E_{S,T}\subseteq Z\) with
\[
K(S,E_{S,T})=0,\qquad K(T,Z\setminus E_{S,T})=0.
\]
Suppose that \(z_S=z_T\). If \(z_S\in E_{S,T}\), monotonicity gives
\[
0<K(S,\{z_S\})\le K(S,E_{S,T})=0.
\]
If \(z_S\notin E_{S,T}\), then \(z_T\in Z\setminus E_{S,T}\), and
\[
0<K(T,\{z_T\})\le K(T,Z\setminus E_{S,T})=0.
\]
Both cases yield a contradiction. Hence the three selected points are distinct. Hence
\[
\kappa(Q)
=\#\{z\in Z:\nu_Q(\{z\})\ne0\}
\ge |A_3|=3.
\]
Together with the attaining three-output staircase channel, this proves exact optimal cardinality three throughout \(U_3\), and completes the result.
% lean: CausalSmith.Stat.LdpAteEfficiencySurface.optimalCardinality_of_mem_R3_of_distinct
\end{enumerate}
\end{proof}
